%% ============================================================
%%  SECCIÓN 2 — Resumen
%%  Responsable: Minilux
%% ============================================================

\section{Resumen}

Se determinó experimentalmente el coeficiente de tensión superficial
del agua mediante dos métodos independientes basados en el equilibrio
estático. En el primer método se empleó una balanza de Mohr--Westphal
con un anillo metálico sumergido en agua pura: la fuerza de tensión
superficial se equilibró mediante jinetillos de masa conocida sobre el
brazo mayor, y el coeficiente $\gamma$ se calculó a partir del equilibrio de torques
y la longitud de contacto del anillo. En el segundo método, se formó
una película jabonosa sobre un dispositivo de tubitos con hilo; al
suspender el sistema, la contracción de la película curva el hilo
inferior, y $\gamma$ se dedujo a partir de las dimensiones geométricas
de la película y el peso del tubo inferior mediante un análisis de
equilibrio de fuerzas.

El primer método, realizado con agua pura, arrojó un coeficiente de
tensión superficial promedio de
$\overline{\gamma}_1 = \qty{39.73 \pm 0.88}{\milli\newton\per\meter}$,
con un error relativo del \qty{45.4}{\percent} respecto al valor
bibliográfico $\gamma_{\text{ref}} = \qty{72.8}{\milli\newton\per\meter}$
\cite{serway}. El segundo método, realizado con solución jabonosa,
produjo $\gamma_2 = \qty{12.21 \pm 1.10}{\milli\newton\per\meter}$,
con un error relativo del \qty{83.2}{\percent}. La marcada diferencia
entre ambos resultados es consistente con la reducción de la tensión
superficial inducida por el surfactante presente en la solución
jabonosa \cite{halliday}.

La discrepancia del primer método respecto al valor bibliográfico se
atribuye principalmente a la imprecisión en la detección del punto de
desprendimiento del anillo, a la posible contaminación del agua y a la
ausencia de un control de temperatura durante el experimento. El segundo
método presenta, además, una alta sensibilidad geométrica: un error de
$\Delta b = \qty{0.1}{\milli\meter}$ en el semiancho mínimo del hilo
se propaga como una variación relativa de $\sim\qty{6}{\percent}$ en
$\gamma_2$, lo que lo hace especialmente vulnerable a errores de
medición.